\section{Introduction}
\label{sec:introduction}

AES acceleration is relevant to communication and data-processing systems that must encrypt
large volumes without excessive latency. CPUs equipped with AES-NI provide a strong reference
for this task, while GPUs offer parallel execution for independent blocks. The practical
question is whether a GPU advantage survives the operations needed to move data through the
accelerator. A fast kernel does not necessarily provide a faster host-to-host processing path.

Existing work demonstrates substantial GPU AES throughput and recognizes the importance of
transfer costs and optimized CPU references. Yet a result measured on one implementation and
timing boundary cannot determine whether offloading is beneficial on another system. Useful
evidence must identify the CPU path, preserve workload pairing, distinguish device computation
from transfer-inclusive processing, and report variability rather than a selected best time.
This is particularly relevant to accessible laptop GPUs, whose suitability cannot be inferred
from results obtained on a different host or accelerator.

We address this problem through a validation-first comparison of AES-NI and a validated CUDA
AES-128 implementation. The study examines ECB and CTR across input sizes and three computing
configurations. It addresses three research questions:

\begin{enumerate}
  \item[\textbf{RQ1}] Under which tested sizes and cipher modes does GPU execution outperform
    AES-NI, and which advantages are consistent across the observed sessions?
  \item[\textbf{RQ2}] How does including transfers change the magnitude and direction of the
    resident-only CPU--GPU comparison?
  \item[\textbf{RQ3}] How do session variability, mode dependence, workload-size effects, and
    differences among systems constrain the interpretation of GPU acceleration?
\end{enumerate}

The paper makes three contributions: a reproducible protocol combining cryptographic
validation, paired workloads and explicit timing scopes; a session-based assessment of
AES-NI and GPU execution on the selected systems; and an interpretation that retains positive,
negative and mixed-session outcomes. The contribution is a defensible comparison of measured
execution paths, not a new AES primitive, a throughput record, or a universal hardware ranking.

Section~\ref{sec:related_work} positions the study within earlier research.
Section~\ref{sec:methodology} defines the platforms, protocol and estimators.
Section~\ref{sec:results} reports the measurements, and Section~\ref{sec:discussion} considers
their practical meaning and limitations. Section~\ref{sec:conclusion} answers the research
questions.
